% !TEX program = xelatex % Указываем, что компилятор - XeLaTeX
\documentclass[aspectratio=169]{beamer}
\usetheme{Madrid}

% после \usetheme{...}
\setbeamersize{text margin left=15pt,text margin right=15pt}


% ставим размер для записей оглавления
\setbeamerfont{section in toc}{size=\scriptsize}
\setbeamerfont{subsection in toc}{size=\scriptsize}
\setbeamerfont{subsubsection in toc}{size=\scriptsize}

% уменьшаем нумерацию/проектированные номера, если они великоваты
\setbeamerfont{section number projected}{size=\scriptsize}
\setbeamerfont{subsection number projected}{size=\scriptsize}

% если пункты выглядят как itemize в ToC
\setbeamerfont{itemize/enumerate body}{size=\scriptsize}


% Уменьшаем отступ для section в оглавлении
\setbeamertemplate{section in toc}[circle] % стандартный кружок
%\setbeamertemplate{section in toc}[square] % можно заменить на квадрат
\setlength{\leftmargini}{0.5em} % уменьшение отступа для itemize в toc
\setlength{\labelsep}{0.5em}    % расстояние между маркером и текстом

\usecolortheme{beaver}

% ---------- Encoding & language (xelatex-friendly) ----------
\usepackage{polyglossia} % Вместо babel для XeLaTeX
\setmainlanguage{russian}
\setotherlanguage{english}

\usepackage{fontspec} % Для работы со шрифтами
\setmainfont{DejaVu Serif}
\setsansfont{DejaVu Sans}
\setmonofont{DejaVu Sans Mono}

% ---------- Common packages ----------
\usepackage{graphicx}
\usepackage{booktabs}
\usepackage{tabularx}
\usepackage{hyperref}
\usepackage{multicol}
\usepackage{ulem}
\usepackage{csquotes}
\usepackage{tikz}
\usetikzlibrary{arrows.meta,positioning,fit,shapes.misc}

% ---------- Listings (Kotlin/Java) ----------
\usepackage{listings} % Используем стандартный listings (inputenc не нужен)
\usepackage[dvipsnames]{xcolor}

\lstdefinelanguage{Kotlin}{
  comment=[l]{//},
  commentstyle={\color{gray}\ttfamily},
  emph={filter, first, firstOrNull, forEach, lazy, map, mapNotNull, println, launch, async, await, withContext, let, run, apply, also, takeIf, takeUnless, repeat, with},
  emphstyle={\color{OrangeRed}},
  identifierstyle=\color{black},
  keywords={!in, !is, abstract, actual, annotation, as, as?, break, by, catch, class, companion, const, constructor, continue, crossinline, data, delegate, do, dynamic, else, enum, expect, external, false, field, file, final, finally, for, fun, get, if, import, in, infix, init, inline, inner, interface, internal, is, lateinit, noinline, null, object, open, operator, out, override, package, param, private, property, protected, public, receiver, reified, return, sealed, set, setparam, super, suspend, tailrec, this, throw, true, try, typealias, typeof, val, var, vararg, when, where, while},
  keywordstyle={\color{NavyBlue}\bfseries},
  morecomment=[s]{/*}{*/},
  morestring=[b]",
  ndkeywords={@Composable, @Preview, @Deprecated, @JvmField, @JvmName, @JvmOverloads, @JvmStatic, @JvmSynthetic, Array, Byte, Double, Float, Int, Integer, Iterable, Long, Runnable, Short, String, Any, Unit, Nothing, Pair, Triple},
  ndkeywordstyle={\color{BurntOrange}\bfseries},
  sensitive=true,
  stringstyle={\color{ForestGreen}\ttfamily},
   keepspaces=true,       % сохраняет пробелы
    columns=flexible       % лучше работает с моноширинными шрифтами
}

\lstdefinelanguage{JavaLite}{
  language=Java,
  basicstyle=\footnotesize\ttfamily,
  keywordstyle=\color{NavyBlue}\bfseries,
  commentstyle=\color{gray}\ttfamily,
  stringstyle=\color{ForestGreen}\ttfamily
}

\lstset{
  extendedchars=true,
  basicstyle=\footnotesize\ttfamily,
  breaklines=true,
  showstringspaces=false,
  frame=single,
  framerule=0.2pt,
  tabsize=2,
  numbers=left,
  numberstyle=\tiny,
  numbersep=6pt
}

% ---------- Title ----------
\title{Объектно-ориентированное программирование в Kotlin}
\subtitle{Лекция}
\author{Александр Глускер}
\institute{Курс «Мобильная разработка»}
\date{\today}




% ---------- TOC at section start ----------
\AtBeginSection[]{
  \begin{frame}{Навигация по лекции}
    \hspace*{12pt}
    { \tiny
    \tableofcontents[currentsection, hideallsubsections]
    }
  \end{frame}
}

\begin{document}

\frame{\titlepage}

%\begin{frame}{Навигация по лекции}
%  \tableofcontents[hideallsubsections]
%\end{frame}

\section{Классы, объекты и конструкторы}

\begin{frame}[fragile]{Классы и объекты}
  \begin{lstlisting}[language=Kotlin]
// Определение класса
class Person(val name: String, var age: Int)

// Создание объекта
val person = Person("Alice", 30)
println(person.name) // Alice
person.age = 31
  \end{lstlisting}
  \begin{itemize}
    \item Класс — шаблон для создания объектов.
    \item \texttt{val} — свойство только для чтения, \texttt{var} — изменяемое.
    \item Первичный конструктор объявляется в заголовке класса.
  \end{itemize}
\end{frame}

\begin{frame}[fragile]{Первичный и вторичный конструкторы}
  \begin{lstlisting}[language=Kotlin]
class Person(val name: String) {
  var age: Int = 0

  // Вторичный конструктор
  constructor(name: String, age: Int) : this(name) {
    this.age = age
  }
}

val p1 = Person("Bob")          // первичный
val p2 = Person("Bob", 25)      // вторичный
  \end{lstlisting}
  \begin{itemize}
    \item Первичный конструктор — часть заголовка класса.
    \item Вторичные конструкторы должны делегировать вызов первичному (\texttt{this(...)}).
  \end{itemize}
\end{frame}

\section{Свойства и функции}

\begin{frame}[fragile]{Свойства с кастомными get/set}
  \begin{lstlisting}[language=Kotlin]
class Temperature {
  var celsius: Double = 0.0
    set(value) {
      field = if (value < -273.15) -273.15 else value
    }
    get() {
      println("Доступ к температуре")
      return field
    }
}
  \end{lstlisting}
  \begin{itemize}
    \item \texttt{field} — резервное поле, автоматически создаётся при наличии кастомного геттера/сеттера.
    \item Геттер/сеттер вызываются при каждом доступе к свойству.
  \end{itemize}
\end{frame}

\section{Вложенные и внутренние классы}

\begin{frame}[fragile]{Nested и Inner классы}
  \begin{lstlisting}[language=Kotlin]
class Outer {
  private val bar = 1

  class Nested {  // nested (статический по умолчанию)
    fun foo() = 2
  }

  inner class Inner {  // inner — имеет доступ к внешнему классу
    fun foo() = bar
  }
}

val nested = Outer.Nested().foo() // OK
val inner = Outer().Inner().foo() // OK
  \end{lstlisting}
\end{frame}

\section{Абстрактные классы и наследование}

\begin{frame}[fragile]{Абстрактные классы и наследование}
  \begin{lstlisting}[language=Kotlin]
abstract class Animal {
  abstract fun makeSound()
  open fun sleep() = println("Zzz")
}

class Dog : Animal() {
  override fun makeSound() = println("Woof!")
  override fun sleep() = println("Dog is sleeping")
}
  \end{lstlisting}
  \begin{itemize}
    \item По умолчанию классы и методы \texttt{final} — нельзя наследовать/переопределять.
    \item Чтобы разрешить — использовать \texttt{open} или \texttt{abstract}.
    \item \texttt{abstract} методы не имеют реализации.
  \end{itemize}
\end{frame}

\section{Companion objects и константы}

\begin{frame}[fragile]{Companion objects}
  \begin{lstlisting}[language=Kotlin]
class MathUtils {
  companion object {
    const val PI = 3.14159
    fun double(x: Int) = x * 2
  }
}

// Использование как статические члены
println(MathUtils.PI)
println(MathUtils.double(5))
  \end{lstlisting}
  \begin{itemize}
    \item \texttt{companion object} — замена статическим методам/полям.
    \item \texttt{const val} — compile-time константа (только примитивы и строки).
  \end{itemize}
\end{frame}

\section{Интерфейсы}

\begin{frame}[fragile]{Интерфейсы}
  \begin{lstlisting}[language=Kotlin]
interface Flyable {
  fun fly()
  fun land() = println("Landed") // default implementation
}

class Bird : Flyable {
  override fun fly() = println("Flying!")
}

val bird = Bird()
bird.fly()
bird.land()
  \end{lstlisting}
  \begin{itemize}
    \item Интерфейсы могут содержать реализацию по умолчанию.
    \item Нет множественного наследования классов, но можно реализовать несколько интерфейсов.
  \end{itemize}
\end{frame}

\section{Области видимости}

\begin{frame}[fragile]{Области видимости}
  \begin{itemize}
    \item \texttt{private} — только в этом классе.
    \item \texttt{protected} — в этом классе и подклассах.
    \item \texttt{internal} — в пределах модуля (аналог package-private в Java).
    \item \texttt{public} — по умолчанию, видно везде.
  \end{itemize}
  \begin{lstlisting}[language=Kotlin]
class Example {
  private fun secret() = "hidden"
  internal fun moduleOnly() = "visible in module"
}
  \end{lstlisting}
\end{frame}

\section{Lateinit и делегированные свойства}

\begin{frame}[fragile]{Lateinit}
  \begin{lstlisting}[language=Kotlin]
class MyActivity {
  lateinit var adapter: RecyclerView.Adapter<*>

  fun onCreate() {
    adapter = MyAdapter()
  }

  fun useAdapter() {
    if (::adapter.isInitialized) {
      // безопасное использование
    }
  }
}
  \end{lstlisting}
  \begin{itemize}
    \item \texttt{lateinit} — откладывает инициализацию \texttt{var}-свойства.
    \item Только для \texttt{var}, не-nullable типов, не примитивов.
    \item Проверка инициализации: \texttt{::prop.isInitialized}.
  \end{itemize}
\end{frame}

\begin{frame}[fragile]{Делегированные свойства}
  \begin{lstlisting}[language=Kotlin]
import kotlin.properties.Delegates

class Example {
  val lazyValue: String by lazy {
    println("Вычисляется один раз!")
    "Hello"
  }

  var observableValue: Int by Delegates.observable(0) { _, old, new ->
    println("Изменено с $old на $new")
  }
}

val e = Example()
println(e.lazyValue) // "Вычисляется один раз!" → "Hello"
e.observableValue = 42 // "Изменено с 0 на 42"
  \end{lstlisting}
\end{frame}

\section{SAM-интерфейсы и extension}

\begin{frame}[fragile]{SAM-интерфейсы}
  \begin{lstlisting}[language=Kotlin]
fun interface OnClickListener {
  fun onClick()
}

// Использование лямбды вместо анонимного класса
button.setOnClickListener { println("Clicked!") }
  \end{lstlisting}
  \begin{itemize}
    \item SAM (Single Abstract Method) — интерфейс с одним абстрактным методом.
    \item Позволяет передавать лямбду вместо реализации.
    \item Объявляется с \texttt{fun interface}.
  \end{itemize}
\end{frame}

\begin{frame}[fragile]{Extension-функции и свойства}
  \begin{lstlisting}[language=Kotlin]
fun String.lastChar(): Char = this[this.length - 1]

val str = "Kotlin"
println(str.lastChar()) // 'n'

// Extension-свойство
val String.isLong: Boolean
  get() = length > 10
  \end{lstlisting}
  \begin{itemize}
    \item Расширяют функциональность существующих типов без наследования.
    \item Компилируются в статические утилитарные методы.
  \end{itemize}
\end{frame}

\section{Data, sealed и enum классы}

\begin{frame}[fragile]{Data классы}
  \begin{lstlisting}[language=Kotlin]
data class User(val name: String, val id: Int)

val user1 = User("Alice", 1)
val user2 = user1.copy(name = "Bob")
println(user1 == user2) // false
println(user1) // User(name=Alice, id=1)
  \end{lstlisting}
  \begin{itemize}
    \item Автоматически генерируют \texttt{equals()}, \texttt{hashCode()}, \texttt{toString()}, \texttt{copy()}.
    \item Все свойства в первичном конструкторе участвуют в сравнении.
  \end{itemize}
\end{frame}

\begin{frame}[fragile]{Sealed классы}
  \begin{lstlisting}[language=Kotlin]
sealed class Result
data class Success(val data: String) : Result()
data class Error(val message: String) : Result()

fun handle(result: Result) = when(result) {
  is Success -> println("Data: ${result.data}")
  is Error -> println("Error: ${result.message}")
}
  \end{lstlisting}
  \begin{itemize}
    \item Ограниченное наследование: подклассы должны быть в том же файле.
    \item Идеальны для моделирования закрытых множеств состояний.
    \item Безопасны в \texttt{when} — компилятор проверяет полноту.
  \end{itemize}
\end{frame}

\begin{frame}[fragile]{Enum классы}
  \begin{lstlisting}[language=Kotlin]
enum class Color(val rgb: Int) {
  RED(0xFF0000),
  GREEN(0x00FF00),
  BLUE(0x0000FF);

  fun printInfo() = println("$name = $rgb")
}

Color.RED.printInfo() // RED = 16711680
for (c in Color.values()) println(c)
  \end{lstlisting}
  \begin{itemize}
    \item Каждая константа — объект enum-класса.
    \item Можно добавлять свойства и методы.
  \end{itemize}
\end{frame}

\begin{frame}{Спасибо за внимание!}
  \centering
  \Large Вопросы?
\end{frame}

\end{document}
